\phantomsection
\addcontentsline{toc}{section}{Abstract}
\markboth{Abstract}{}

\section*{Abstract}

This dissertation develops a three-stage statistical framework for predicting health insurance costs.  The analysis uses a publicly available benchmark dataset containing one million synthetic policyholder records.  Because the dataset has a clean, simulated structure, the predictive accuracy reported here is substantially higher than would be expected on real-world insurance data, and the specific performance figures should not be generalised beyond this benchmark setting.  The methodology, however, is designed to be transferable to other datasets and domains.

In the first stage, a deep neural network is trained to predict insurance charges from policyholder characteristics.  Its accuracy is compared against three established statistical models: a traditional insurance pricing model, a random forest and a gradient-boosted tree ensemble.  A five-phase selection process evaluates different network configurations, and the final model is assessed on held-out data to ensure the results reflect genuine predictive ability rather than memorisation.

In the second stage, computer simulation generates a large population of realistic but synthetic policyholder profiles based on the statistical patterns observed in the data.  The trained network then predicts costs for each simulated profile, producing a full distribution of predicted costs rather than predictions at individual data points alone.  This distribution is checked for consistency with the original data through a suite of validation diagnostics.  Because the network achieves a near-perfect fit and the simulation draws from the same statistical patterns used in training, the simulated distribution closely mirrors the original charge distribution; the primary purpose of this stage is therefore to define the high-cost regions used in the next stage.

In the third stage, each predicted cost is broken down into the individual contributions of policyholder characteristics such as age, smoking status, coverage level and medical history.  This decomposition is performed both across the full simulated population and specifically among the highest-cost profiles, using a method rooted in cooperative game theory that fairly allocates each prediction among the contributing variables.

On this benchmark dataset, the deep neural network captures over 99\% of the variation in insurance costs, outperforming all three comparison models.  The contribution analysis identifies smoking status, coverage level and medical history as the strongest overall cost drivers.  Among the highest-cost profiles, however, the ranking changes: coverage level, medical history and family medical history become more important than smoking status.  Part of this shift is mechanical, because smoking status is near-universal among the highest-cost profiles and therefore loses its ability to distinguish between them; part reflects the genuine explanatory role of the remaining variables within the upper tail of the cost distribution.  This finding demonstrates that the factors driving extreme predicted costs are not necessarily the same as those most important on average.

All conclusions are internal to the fitted model and this benchmark dataset.  No causal claims are made, and the modelling assumptions are stated explicitly throughout the dissertation.

\vspace{4mm}

\noindent\textbf{Keywords:} deep neural network, health insurance cost prediction, Monte Carlo simulation, Shapley values, explainable machine learning, generalised linear model, random forest, XGBoost, statistical learning.
